% !TEX program = xelatex
\documentclass[11pt,a4paper]{ctexart}
% 零基础讲义公共导言区
\usepackage{amsmath,amssymb,amsthm,mathtools}
\usepackage{geometry}
\geometry{left=2.0cm,right=2.0cm,top=2.2cm,bottom=2.0cm}
\usepackage{booktabs}
\usepackage{array}
\usepackage{enumitem}
\usepackage{tcolorbox}
\tcbuselibrary{skins,breakable}
\usepackage{xcolor}
\usepackage{fancyhdr}
\usepackage{titlesec}
\usepackage{multirow}
\usepackage{longtable}

% ---------- 颜色 ----------
\definecolor{mainblue}{RGB}{25,70,130}
\definecolor{acc1}{RGB}{176,58,46}      % 定义/公式
\definecolor{acc2}{RGB}{20,105,85}      % 直觉/类比
\definecolor{acc3}{RGB}{120,80,150}     % 你可能想问
\definecolor{acc4}{RGB}{160,110,20}     % 手把手算
\definecolor{acc5}{RGB}{40,90,160}      % 考试怎么考
\definecolor{soft}{RGB}{246,247,250}

% ---------- 数学宏 ----------
\newcommand{\Hx}[1]{H\!\left(#1\right)}
\newcommand{\Hc}[1]{H_{c}\!\left(#1\right)}
\newcommand{\I}[1]{I\!\left(#1\right)}
\newcommand{\lb}{\log_{2}}
\DeclareMathOperator{\E}{\mathbb{E}}
\newcommand{\bit}{\,\text{bit}}
\newcommand{\p}[1]{p\!\left(#1\right)}

% ---------- 五种盒子 ----------
\newtcolorbox{defbox}[1]{enhanced,breakable,
  colback=soft,colframe=acc1,boxrule=0.7pt,arc=2pt,
  left=6pt,right=6pt,top=5pt,bottom=5pt,
  title={\bfseries #1},fonttitle=\color{white},colbacktitle=acc1,
  attach boxed title to top left={xshift=6pt,yshift=-2.6pt},
  boxed title style={arc=1pt,boxrule=0pt}}

\newtcolorbox{idea}[1]{enhanced,breakable,
  colback=white,colframe=acc2,boxrule=0.7pt,arc=2pt,
  left=6pt,right=6pt,top=5pt,bottom=5pt,
  title={\bfseries #1},fonttitle=\color{white},colbacktitle=acc2,
  attach boxed title to top left={xshift=6pt,yshift=-2.6pt},
  boxed title style={arc=1pt,boxrule=0pt}}

\newtcolorbox{qa}[1]{enhanced,breakable,
  colback=white,colframe=acc3,boxrule=0.7pt,arc=2pt,
  left=6pt,right=6pt,top=5pt,bottom=5pt,
  title={\bfseries #1},fonttitle=\color{white},colbacktitle=acc3,
  attach boxed title to top left={xshift=6pt,yshift=-2.6pt},
  boxed title style={arc=1pt,boxrule=0pt}}

\newtcolorbox{calc}[1]{enhanced,breakable,
  colback=white,colframe=acc4,boxrule=0.7pt,arc=2pt,
  left=6pt,right=6pt,top=5pt,bottom=5pt,
  title={\bfseries #1},fonttitle=\color{white},colbacktitle=acc4,
  attach boxed title to top left={xshift=6pt,yshift=-2.6pt},
  boxed title style={arc=1pt,boxrule=0pt}}

\newtcolorbox{exam}[1]{enhanced,breakable,
  colback=white,colframe=acc5,boxrule=0.7pt,arc=2pt,
  left=6pt,right=6pt,top=5pt,bottom=5pt,
  title={\bfseries #1},fonttitle=\color{white},colbacktitle=acc5,
  attach boxed title to top left={xshift=6pt,yshift=-2.6pt},
  boxed title style={arc=1pt,boxrule=0pt}}

% 兼容旧环境名
\newenvironment{thmbox}[1]{\begin{defbox}{#1}}{\end{defbox}}
\newenvironment{keypoint}[1]{\begin{idea}{#1}}{\end{idea}}
\newenvironment{warning}[1]{\begin{qa}{#1}}{\end{qa}}
\newenvironment{example}[1]{\begin{calc}{#1}}{\end{calc}}
\newenvironment{problem}[1]{\begin{defbox}{#1}}{\end{defbox}}

% ---------- 标题格式 ----------
\titleformat{\section}{\Large\bfseries\color{acc1}}{第\,\thesection\,讲}{0.6em}{}
\titleformat{\subsection}{\large\bfseries\color{mainblue}}{\thesubsection}{0.5em}{}
\titleformat{\subsubsection}{\normalsize\bfseries\color{acc2}}{\thesubsubsection}{0.4em}{}

% ---------- 页眉页脚 ----------
\pagestyle{fancy}
\fancyhf{}
\fancyhead[L]{\small 信息论基础\ \textbar\ 零基础讲义}
\fancyhead[R]{\small\leftmark}
\fancyfoot[C]{\small\thepage}
\renewcommand{\headrulewidth}{0.4pt}

\allowdisplaybreaks
\setlength{\parskip}{2.5pt}
\linespread{1.08}

% ---------- 特殊符号兜底（避免缺字）----------
\usepackage{newunicodechar}
\newunicodechar{✓}{\ensuremath{\checkmark}}
\newunicodechar{✗}{\ensuremath{\times}}
\xeCJKDeclareCharClass{CJK}{"2460 -> "24FF}   % ① ② ③ 等圈号交给中文字体
\newunicodechar{→}{\ensuremath{\rightarrow}}
\newunicodechar{≥}{\ensuremath{\ge}}
\newunicodechar{≤}{\ensuremath{\le}}
\newunicodechar{≠}{\ensuremath{\ne}}
\newunicodechar{×}{\ensuremath{\times}}
\newunicodechar{≈}{\ensuremath{\approx}}
\newunicodechar{∞}{\ensuremath{\infty}}
\newunicodechar{∈}{\ensuremath{\in}}
\newunicodechar{⊕}{\ensuremath{\oplus}}
\newunicodechar{∑}{\ensuremath{\sum}}
\newunicodechar{∫}{\ensuremath{\int}}
\newunicodechar{√}{\ensuremath{\surd}}
\newunicodechar{σ}{\ensuremath{\sigma}}
\newunicodechar{π}{\ensuremath{\pi}}
\newunicodechar{λ}{\ensuremath{\lambda}}
\newunicodechar{μ}{\ensuremath{\mu}}
\newunicodechar{Δ}{\ensuremath{\Delta}}
\newunicodechar{α}{\ensuremath{\alpha}}
\newunicodechar{β}{\ensuremath{\beta}}
\newunicodechar{ε}{\ensuremath{\varepsilon}}

\linespread{1.05}
\setlength{\parskip}{2pt}

\begin{document}

\begin{center}
{\LARGE\bfseries 第 3 章 保命速查（互信息）}\\[6pt]
{\large 只看这一份，就能看懂第二次作业的四道题}\\[6pt]
{\small 全章你只需要记住：\textbf{一条定义} + \textbf{三个工具}。其他全是它们的变形。}
\end{center}

\vspace{-2pt}
\begin{idea}{先给你吃颗定心丸：第 3 章其实只有一句话}
\[
\boxed{\ \I{X;Y}=\Hx{X}-\Hx{X|Y}\ }
\qquad\text{人话：知道 $Y$ 之后，我对 $X$ 的不确定性减少了多少。}
\]
减少的那部分，就是 $Y$ 里关于 $X$ 的“信息量”，所以叫\textbf{互}信息。

\textbf{为什么这一章看起来吓人}：因为同一个数被写成了 3 种样子，书上不停换着用，你以为有 3 个新概念——其实是一个东西的三个马甲。
\end{idea}

\section{一个数，三个马甲（背这三行就够）}

\begin{defbox}{互信息的三种等价写法}
\begin{align*}
\I{X;Y}&=\Hx{X}-\Hx{X|Y} &&\text{（原来的不确定性 减 知道 $Y$ 后剩下的）}\\
&=\Hx{Y}-\Hx{Y|X} &&\text{（换个方向算，结果一样，这叫对称性）}\\
&=\Hx{X}+\Hx{Y}-\Hx{X,Y} &&\text{（分别算再加，减去重复算的）}
\end{align*}
\textbf{考试 90\% 的题就是在这三条里挑一条代入。}看到“求 $\I{X;Y}$”，先看哪个量已知，就用对应的那条。
\end{defbox}

\begin{qa}{三条为什么相等？只需两步}
第一步（链式法则，第 2 章学的）：
\[
\Hx{X,Y}=\Hx{X}+\Hx{Y|X}.
\]
第二步（移项）：
\[
\Hx{Y|X}=\Hx{X,Y}-\Hx{X}\ \Longrightarrow\ \Hx{X}-\Hx{X|Y}=\Hx{Y}-\Hx{Y|X}.
\]
把左边的 $\Hx{X|Y}=\Hx{X,Y}-\Hx{Y}$ 代进第一条定义，就得到第三条：
$\I{X;Y}=\Hx{X}-(\Hx{X,Y}-\Hx{Y})=\Hx{X}+\Hx{Y}-\Hx{X,Y}$。\quad 就这么简单，没有更多东西了。
\end{qa}

\section{手算一遍（作业 2.7 那张表，6 步出结果）}

\begin{calc}{把数字代进去，你就彻底懂了}
联合分布表：
\begin{center}
\renewcommand{\arraystretch}{1.25}
\begin{tabular}{c|cc|c}
\toprule
 & $Y=0$ & $Y=1$ & 行和\\
\midrule
$X=0$ & $1/3$ & $1/3$ & $2/3$\\
$X=1$ & $1/3$ & $0$ & $1/3$\\
\midrule
列和 & $2/3$ & $1/3$ & $1$\\
\bottomrule
\end{tabular}
\end{center}

\textbf{第 1 步}：算两个单变量熵（第 2 章的老活）
\[
\Hx{X}=-\tfrac23\lb\tfrac23-\tfrac13\lb\tfrac13=0.667\times0.585+0.333\times1.585=0.918\bit,\qquad \Hx{Y}=0.918\bit .
\]

\textbf{第 2 步}：算联合熵（表里 3 个非零格子且都等于 $\frac13$，直接 $\lb3$）
\[
\Hx{X,Y}=\lb 3=1.585\bit .
\]

\textbf{第 3 步}：用链式法则拿条件熵（不用重新算！）
\[
\Hx{X|Y}=\Hx{X,Y}-\Hx{Y}=1.585-0.918=0.667\bit .
\]

\textbf{第 4 步}：套第一条马甲
\[
\I{X;Y}=\Hx{X}-\Hx{X|Y}=0.918-0.667=0.251\bit .
\]

\textbf{第 5 步}：换第二条马甲验算
$\Hx{Y|X}=1.585-0.918=0.667$，于是 $\I{X;Y}=0.918-0.667=0.251\bit$ ✓ 一样。

\textbf{第 6 步}：换第三条马甲再验算
$\Hx{X}+\Hx{Y}-\Hx{X,Y}=0.918+0.918-1.585=0.251\bit$ ✓ 还是一样。

\textbf{体会}：三条马甲算出同一个数，所以考试时\textbf{随便挑最省事的那条}，其他两条拿来验算。这就是第 3 章全部的计算题。
\end{calc}

\section{证明题只用这三个工具（背下来就够用）}

\begin{defbox}{工具 1：链式法则（拆开的唯一办法）}
\[
\Hx{X,Y}=\Hx{X}+\Hx{Y|X},\qquad
\I{X;YZ}=\I{X;Y}+\I{X;Z|Y}.
\]
人话：想一起处理两个东西，就“先处理一个，再处理另一个”。
\end{defbox}

\begin{defbox}{工具 2：条件不增熵（把“条件”摘掉）}
\[
\Hx{X|Y}\le\Hx{X},\qquad \Hx{X|Y,Z}\le\Hx{X|Y}.
\]
人话：知道得越多，不确定性只会更小。等号成立当且仅当独立。
\end{defbox}

\begin{defbox}{工具 3：互信息非负（把“互信息”扔掉）}
\[
\I{X;Y}\ge 0,\qquad \I{X;Y|Z}\ge0 .
\]
人话：知道点东西总比不知道强，最差也是白知道。
\end{defbox}

\begin{calc}{综合演示：证明 $\I{X;Y}\le\Hx{X}$（书上的“性质”，也是作业 3.7 第(1)问的双胞胎）}
一行就完：
\[
\I{X;Y}\overset{\text{马甲1}}{=}\Hx{X}-\underbrace{\Hx{X|Y}}_{\ge0}\ \le\ \Hx{X}.
\]
（因为条件熵一定 $\ge 0$。）\\
\textbf{看出来了吗}：所谓“证明题”，就是\textbf{选一个马甲 + 用一次工具}。全书第 3 章的证明题都是这个套路。
\end{calc}

\section{四道作业题，各只需要哪两行}

\begin{center}
\renewcommand{\arraystretch}{1.35}
\small
\begin{tabular}{p{1.5cm}|p{5.2cm}|p{6.4cm}}
\toprule
题号 & 它要你干什么 & 你只需要用到 \\
\midrule
3.2 & 证 $\I{X;Y}\ge\I{X;Y|Z}$ & 关键恒等式 $\I{X;Y}-\I{X;Y|Z}=\I{X;Z}-\I{X;Z|Y}$\\
    & （教材漏了条件，实际要 $X\to Y\to Z$） & 然后：若 $\I{X;Z|Y}=0$，则差 $=\I{X;Z}\ge0$ \\
\midrule
3.3 & 证 $\Hx{X,Y,Z}=\Hx{X,Z}+\Hx{Y|X}-\I{Y;Z|X}$ & 工具 1 拆一次 + 把 $\I{Y;Z|X}$ 的定义挪项 \\
\midrule
3.7 & 给 $Z=X+Y$、$X\perp Y$，证八个式子 & 只用到两件事：\\
    & & ① $Z$ 是 $X,Y$ 的函数 $\Rightarrow\Hx{Z|X,Y}=0$\\
    & & ② $X\perp Y\Rightarrow\Hx{X|Y}=\Hx{X}$\\
\midrule
3.10 & 数据处理：$X\to Y\to Z$，证 $\I{X;Z}\le\I{X;Y}$ & 工具 3：$\I{X;Z}=\I{X;Y}-\underbrace{\I{X;Y|Z}}_{\ge0}$ \\
\bottomrule
\end{tabular}
\end{center}

\begin{exam}{如果你只想“把作业交上去”，最低限度写这些}
\begin{enumerate}[leftmargin=1.6em,itemsep=2pt,topsep=2pt]
  \item \textbf{3.2}：先写“教材原题缺条件”，再写恒等式那两行，再补一句“当 $\I{X;Z|Y}=0$（即 $X\to Y\to Z$）时成立”，并给出反例 $X,Y$ 独立、$Z=X\oplus Y$ 时 $\I{X;Y}=0<1=\I{X;Y|Z}$。
  \item \textbf{3.3}：写 $\Hx{X,Y,Z}=\Hx{X,Z}+\Hx{Y|X,Z}$（工具 1），再把 $\Hx{Y|X,Z}=\Hx{Y|X}-\I{Y;Z|X}$ 代入即得。
  \item \textbf{3.7}：八问里至少会做 (5)(6)：$\I{X,Y;Z}=\Hx{Z}-\Hx{Z|X,Y}=\Hx{Z}-0=\Hx{Z}$，
        $\I{X;Y,Z}=\Hx{X}-\Hx{X|Y,Z}=\Hx{X}-0=\Hx{X}$（因为 $Z$ 由 $X,Y$ 决定，故 $\Hx{Z|X,Y}=0$）。
        剩下每问都是同两个工具。
  \item \textbf{3.10}：$\I{X;Z}=\I{X;Y,Z}-\I{X;Y|Z}=\I{X;Y}+\I{X;Z|Y}-\I{X;Y|Z}$，而 $X\to Y\to Z$ 使 $\I{X;Z|Y}=0$，故 $\I{X;Z}=\I{X;Y}-\I{X;Y|Z}\le\I{X;Y}$。
\end{enumerate}
这些都能在 `作业2-第3章.pdf` 里找到逐步展开的版本，照着抄思路即可。
\end{exam}

\section{能不能跳过？给你一句实话}

\begin{qa}{问题：第 3 章能不能不学？}
\textbf{作业可以少做，概念不能跳过——但也真的不需要全懂。}

\textbf{不学它会付出什么代价（按考试权重排）}：
\begin{enumerate}[leftmargin=1.6em,itemsep=1pt,topsep=2pt]
  \item \textbf{第 5 章的信道容量全部建立在这一章上}：$C=\max_{p(x)}\I{X;Y}$。不会 $\I{X;Y}$，第 5 章的容量题（很可能有大题）直接归零；
  \item \textbf{第 4 章的编码定理}里 $\bar L\ge\Hx{X}$ 的论证也要用到互信息；
  \item 但——\textbf{只要求你会算 $\I{X;Y}$ 和会用那三个工具}，也就是这一份速查的全部内容。
\end{enumerate}

\textbf{最小补救方案（按性价比排序，1 小时能搞定）}：
\begin{enumerate}[leftmargin=1.6em,itemsep=1pt,topsep=2pt]
  \item 先把第 2 节那个 2×2 表的手算，自己动笔算一遍（20 分钟，必须动手）；
  \item 把“三个马甲”背下来，能默写（10 分钟）；
  \item 把“三个工具”背下来（10 分钟）；
  \item 拿 `作业2-第3章.pdf` 当答案册，对着看 3.10 和 3.7(5)(6) 怎么用工具（20 分钟）。
\end{enumerate}
\end{qa}

\begin{idea}{一句话总结这份文档}
\textbf{互信息 = $\Hx{X}-\Hx{X|Y}$；证明题 = 换个马甲 + 用一次“条件不增熵/非负性”；计算题 = 先算三个熵，再套马甲。}就这些。
\end{idea}

\end{document}
